\chapter*{German--English Glossary}
\addcontentsline{toc}{chapter}{German--English Glossary}
\markboth{German--English Glossary}{German--English Glossary}

The mathematics in this book is written in English, while business and
economics degrees across the German-speaking countries set their
examinations in German. The terms below are grouped by subject
rather than alphabetically, because that is how they are revised. Where two
German terms are in common use for one object, both are given.

\section*{Functions and curves}

\begin{itemize}[leftmargin=1.4em,itemsep=0.15ex]
  \item \de{Definitionsbereich} --- domain; \de{Wertebereich},
        \de{Bildmenge} --- range, image
  \item \de{Nullstelle} --- root, zero; \de{Polstelle} --- pole;
        \de{Sprungstelle} --- jump discontinuity
  \item \de{Asymptote} --- asymptote; \de{schiefe Asymptote} --- oblique
        asymptote
  \item \de{gerade} / \de{ungerade Funktion} --- even / odd function
  \item \de{Umkehrfunktion} --- inverse function; \de{Verkettung} ---
        composition
  \item \de{Verschiebung}, \de{Streckung}, \de{Spiegelung} --- shift,
        stretch, reflection
  \item \de{Betrag} --- absolute value; \de{Ungleichung} --- inequality
  \item \de{Potenz}, \de{Wurzel}, \de{Logarithmus} --- power, root,
        logarithm; \de{Winkelfunktion} --- trigonometric function
\end{itemize}

\section*{Limits, derivatives, curve discussion}

\begin{itemize}[leftmargin=1.4em,itemsep=0.15ex]
  \item \de{Folge} --- sequence; \de{Reihe} --- series;
        \de{Grenzwert} --- limit; \de{Stetigkeit} --- continuity
  \item \de{Ableitung} --- derivative; \de{Produktregel},
        \de{Quotientenregel}, \de{Kettenregel} --- product, quotient, chain
        rule
  \item \de{Monotonie} --- monotonicity; \de{streng monoton steigend} ---
        strictly increasing
  \item \de{Kr\"ummung} --- curvature; \de{konvex} / \de{konkav} ---
        convex / concave; \de{Wendepunkt} --- inflection point
  \item \de{Hochpunkt}, \de{Tiefpunkt} --- local maximum, local minimum;
        \de{Extremum} --- extremum; \de{kritischer Punkt} --- critical
        point
  \item \de{Kurvendiskussion} --- (complete) curve discussion, curve
        sketching
  \item \de{Taylor-Approximation} --- Taylor approximation;
        \de{Restglied} --- remainder term
\end{itemize}

\section*{Several variables}

\begin{itemize}[leftmargin=1.4em,itemsep=0.15ex]
  \item \de{partielle Ableitung} --- partial derivative;
        \de{Richtungsableitung} --- directional derivative
  \item \de{totales Differential} --- total differential;
        \de{Gradient} --- gradient
  \item \de{Hesse-Matrix} --- Hessian; \de{Jacobi-Matrix} --- Jacobian
  \item \de{Niveaulinie}, \de{Niveaumenge} --- level curve, level set
  \item \de{Subniveaumenge}, \de{untere Konturmenge} $W(c)$ --- sublevel
        set, lower contour set $L_{c}$
  \item \de{Superniveaumenge}, \de{obere Konturmenge} $B(c)$ ---
        superlevel set, upper contour set $U_{c}$
  \item \de{Sattelpunkt} --- saddle point; \de{quasikonkav} ---
        quasiconcave
  \item \de{homogen vom Grad} --- homogeneous of degree; \de{Satz \"uber
        implizite Funktionen} --- implicit function theorem
  \item \de{offen}, \de{abgeschlossen}, \de{beschr\"ankt},
        \de{kompakt} --- open, closed, bounded, compact
\end{itemize}

\section*{Optimization}

\begin{itemize}[leftmargin=1.4em,itemsep=0.15ex]
  \item \de{Zielfunktion} --- objective function;
        \de{Nebenbedingung}, \de{Restriktion} --- constraint
  \item \de{zul\"assiger Bereich} --- feasible region;
        \de{bindend} / \de{nicht bindend} --- binding / slack
  \item \de{Lagrange-Multiplikator} --- Lagrange multiplier;
        \de{Schattenpreis} --- shadow price
  \item \de{ger\"anderte Hesse-Matrix} --- bordered Hessian
  \item \de{Kuhn-Tucker-Bedingungen} (KTB) --- Kuhn--Tucker conditions;
        \de{komplement\"are Schlupfbedingung} --- complementary slackness
  \item \de{Envelope-Theorem}, \de{H\"ullkurvensatz} --- envelope theorem;
        \de{komparative Statik} --- comparative statics
  \item \de{Substitutionsmethode} --- substitution method;
        \de{konvexes Programm} --- convex programme
\end{itemize}

\section*{Linear algebra and linear programming}

\begin{itemize}[leftmargin=1.4em,itemsep=0.15ex]
  \item \de{lineares Gleichungssystem} --- system of linear equations;
        \de{Gau\ss-Elimination} --- Gaussian elimination
  \item \de{Zeilenstufenform} --- row-echelon form;
        \de{Pivot} --- pivot; \de{freie Variable} --- free variable
  \item \de{Determinante} --- determinant; \de{Rang} --- rank;
        \de{Inverse} --- inverse
  \item \de{Eigenwert}, \de{Eigenvektor} --- eigenvalue, eigenvector;
        \de{Definitheit} --- definiteness
  \item \de{Basisl\"osung} --- basic solution;
        \de{Simplexverfahren} --- simplex method
  \item \de{Dualit\"at} --- duality; \de{Endtableau} --- final tableau;
        \de{Sensitivit\"atsanalyse} --- sensitivity analysis
\end{itemize}

\section*{Integration and financial mathematics}

\begin{itemize}[leftmargin=1.4em,itemsep=0.15ex]
  \item \de{Stammfunktion} --- antiderivative;
        \de{unbestimmtes} / \de{bestimmtes Integral} --- indefinite /
        definite integral
  \item \de{Substitution}, \de{partielle Integration} --- substitution,
        integration by parts
  \item \de{uneigentliches Integral} --- improper integral
  \item \de{Barwert} --- present value; \de{Endwert} --- future value
  \item \de{Rente}, \de{Annuit\"at} --- annuity; \de{ewige Rente} ---
        perpetuity
  \item \de{Zinssatz} --- interest rate; \de{stetige Verzinsung} ---
        continuous compounding; \de{Abzinsung} --- discounting
\end{itemize}

\section*{Economic quantities}

\begin{itemize}[leftmargin=1.4em,itemsep=0.15ex]
  \item \de{Grenzkosten} --- marginal cost; \de{Grenznutzen} --- marginal
        utility; \de{Grenzprodukt} --- marginal product
  \item \de{Durchschnittskosten} --- average cost; \de{Fixkosten} ---
        fixed cost
  \item \de{Elastizit\"at} --- elasticity
  \item \de{Isoquante} --- isoquant; \de{Indifferenzkurve} ---
        indifference curve; \de{Budgetgerade} --- budget line
  \item \de{Grenzrate der Substitution} --- marginal rate of substitution
        (MRS); \de{Grenzrate der technischen Substitution} --- MRTS
  \item \de{Angebot}, \de{Nachfrage} --- supply, demand;
        \de{Konsumentenrente} --- consumer surplus
  \item \de{Skalenertr\"age} --- returns to scale;
        \de{Produktionsfunktion} --- production function
\end{itemize}
